\begin{table}[!tbp]\centering
\caption{Robustness to the Venezuelan immigration wave}\label{tab:migration}
\begin{threeparttable}
\footnotesize
\setlength{\tabcolsep}{4pt}
\begin{tabular}{lcccc}
\toprule
Specification & Log hourly wage & Employment & Formal empl. & Self-employment \\
\midrule
Baseline & -0.021 & -0.025$^{***}$ & -0.046$^{***}$ & 0.036$^{***}$ \\
 & (0.012) & (0.006) & (0.007) & (0.006) \\
Department$\times$quarter FE & -0.020 & -0.015$^{**}$ & -0.041$^{***}$ & 0.034$^{***}$ \\
 & (0.013) & (0.006) & (0.008) & (0.006) \\
Excluding Lima and Callao & -0.021 & -0.018$^{*}$ & -0.042$^{***}$ & 0.037$^{***}$ \\
 & (0.022) & (0.009) & (0.010) & (0.010) \\
Excluding top-8 migrant-hosting departments & -0.051 & -0.012 & -0.049$^{***}$ & 0.045$^{**}$ \\
 & (0.032) & (0.013) & (0.015) & (0.016) \\
\bottomrule
\end{tabular}
\begin{tablenotes}\footnotesize
\item Notes: $\mathrm{Low}\times\mathrm{Post}$ DiD estimates under specifications that address the Venezuelan immigration inflow. Department$\times$quarter fixed effects absorb any local shock common to both education groups within a department-quarter, including migrant arrivals and regularization waves. The top-8 exclusion drops Lima, Callao, La Libertad, Arequipa, Lambayeque, Piura, Ica, and Tumbes, the departments hosting the large majority of the Venezuelan population. Department-clustered standard errors; ENAHO survey weights. $^{*}p<0.1$, $^{**}p<0.05$, $^{***}p<0.01$.
\end{tablenotes}
\end{threeparttable}\end{table}
